\BourbakiExercises{Exercises}

\begin{enumerate}
\def\labelenumi{\arabic{enumi})}
\tightlist
\item
  Let \(\mathrm{K}\) be a commutative field and \(\mathrm{A}\) a K-algebra. We call an A-module M absolutely without radical if for every extension \(\mathrm{L}\) of \(\mathrm{K}\) , the \(\mathrm{A}_{(\mathrm{L})}\) -module \(\mathrm{M}_{(\mathrm{L})}\) is without radical. We say that the algebra \(\mathrm{A}\) is absolutely without radical if the A-module \(\mathrm{A}_s\) is absolutely without radical.
\end{enumerate}

\begin{enumerate}
\def\labelenumi{\alph{enumi})}
\item
  A commutative algebra that is absolutely without radical is separable (V, §15, No.~2, p.~119, Definition 1).
\item
  Every submodule of a module that is absolutely without radical is absolutely without radical; a direct sum of modules that are absolutely without radical is absolutely without radical.
\item
  Let M be a finitely generated A-module; then M is absolutely without radical if and only if there exists a perfect field P that is an algebraic extension of K such that \(\mathfrak{R}_{\mathrm{A}_{(\mathrm{P})}}(\mathrm{M}_{(\mathrm{P})}) = 0\) .
\item
  A semisimple A-module M is absolutely without radical if and only if every simple module belonging to the support of M is.
\item
  Let S be a simple A-module and D its commutant. Prove that the following properties are equivalent:
\end{enumerate}

\begin{enumerate}
\def\labelenumi{(\roman{enumi})}
\item
  The A-module S is absolutely without radical.
\item
  The K-algebra \(\mathrm{D}\) is absolutely without radical.
\item
  The center of \(\mathrm{D}\) is a separable extension of \(\mathrm{K}\) .
\end{enumerate}

\(f\) ) A semisimple algebra is absolutely without radical if and only if the center of each of its simple components is a separable extension of K.

\(g\) ) Prove that an A-module that is absolutely without radical and finite-dimensional over \(\mathbf{K}\) is absolutely semisimple.

\begin{enumerate}
\def\labelenumi{\arabic{enumi})}
\setcounter{enumi}{1}
\tightlist
\item
  Let \(A_{1}\) and \(A_{2}\) be K-algebras, A the algebra \(A_{1} \otimes_{K} A_{2}\) , \(M_{i}\) an \(A_{i}\) -module (for i = 1, 2), and M the A-module \(M_{1} \otimes_{K} M_{2}\) .
\end{enumerate}

\begin{enumerate}
\def\labelenumi{\alph{enumi})}
\item
  Suppose that the \(A_{1}\) -module \(M_{1}\) is absolutely without radical. Prove the inclusion \(\mathfrak{R}_{\mathrm{A}}(\mathrm{M}) \subset \mathrm{M}_{1} \otimes_{\mathrm{K}} \mathfrak{R}_{\mathrm{A}_{2}}(\mathrm{M}_{2})\) .
\item
  If the modules \(M_{1}\) and \(M_{2}\) are absolutely without radical, then so is the A-module M.
\end{enumerate}

\begin{enumerate}
\def\labelenumi{\arabic{enumi})}
\setcounter{enumi}{2}
\item
  Let K be a commutative field, and let A and B be simple K-algebras. Suppose that A is of finite rank over its center Z, that Z is an algebraic extension of K, and that the K-algebra B is absolutely without radical (Exercise 1). Prove that the ring \(A \otimes_{K} B\) is absolutely flat (VIII, p.~171, Exercise 27; observe that every element of \(A \otimes_{K} B\) belongs to a subalgebra of the form \(A_{1} \otimes_{K} B\) , where \(A_{1}\) is a simple algebra of finite rank over K).
\item
  Let \(\mathrm{K}\) be a commutative field of characteristic \(p > 0\) , and let \(a\) be an element of \(\mathrm{K} - \mathrm{K}^p\) . Denote the (finite-dimensional) K-algebra \(\mathrm{K}[\mathrm{X}] / ((\mathrm{X}^p - a)^2)\) by A. Prove that the radical \(\mathfrak{r}\) of A has square zero, that the K-algebra \(\mathrm{A} / \mathfrak{r}\) is isomorphic to the radical extension \(\mathrm{L} = \mathrm{K}[\mathrm{X}] / (\mathrm{X}^{p} - a)\) of K, and that A does not contain any K-subalgebras isomorphic to L.
\item
  Let \(\mathrm{K}\) be a commutative field and \(\mathrm{L}\) a finite set. Denote the algebra constructed in Exercise 17 of VIII, p.~169 by B, the central element \(1 - \sum_{\lambda} c_{\lambda} + \sum_{\lambda \neq \mu} r_{\lambda \mu}\) of \(\mathrm{B}\) by \(\varepsilon\) , and the algebra \(\mathrm{B} / \mathrm{B}\varepsilon\) by A. The K-algebra A is finite-dimensional, its radical \(\mathfrak{r}\) has square zero, and \(\mathrm{A} / \mathfrak{r}\) is isomorphic to \(\mathrm{K}^{\mathrm{L}}\) . Prove that A is not isomorphic to a product of two nonzero algebras.
\end{enumerate}

¶ 6) Let K be a commutative field and E an extension of K. Prove that the following properties are equivalent:

\begin{enumerate}
\def\labelenumi{(\roman{enumi})}
\item
  For every separable extension L of K, the ring \(E_{(L)}\) is semisimple.
\item
  The field \(\mathrm{E}\) is a purely inseparable extension of a separable extension of \(\mathrm{K}\) of finite degree.
\end{enumerate}

(Assuming (i), use the corollary of V, §17, No.~2, p.~140 and Exercise 4 of V, §2, p.~146 to prove that E is an algebraic extension of K. Then, prove that the separable closure \(E_{s}\) of K in E has finite degree. To do this, note that in the opposite case, the ring \(\mathrm{E}_{(\mathrm{L})}\) , where L is a separable closure of K, would contain families of orthogonal idempotents \((e_{1},\ldots,e_{n})\) with n arbitrarily large. Assuming (ii), prove that if \(\mathrm{E}_{s(\mathrm{L})}\) is isomorphic to \(\prod_{j=1}^{r}C_{j}\) , then \(\mathrm{E}_{(\mathrm{L})}\) is isomorphic to \(\prod_{j=1}^{r}D_{j}\) , where \(D_{j}\) is a purely inseparable extension of \(C_{j}\) .)

¶ 7) Let E be a p-radical extension of a commutative field K of characteristic p \textgreater{} 0. Suppose that the ring \(E \otimes_{K} E\) is Noetherian.

\begin{enumerate}
\def\labelenumi{\alph{enumi})}
\item
  Prove that there exists an integer \(k\) such that \(\mathrm{E} \subset \mathrm{K}^{p^{-k}}\) (reason by contradiction: construct a sequence \((x_{n})\) in \(\mathrm{E}\) such that \((x_{n} \otimes 1 - 1 \otimes x_{n})^{p^{n-1}} \neq 0\) and \((x_{n} \otimes 1 - 1 \otimes x_{n})^{p^n} = 0\) ).
\item
  Let \(\mathrm{F} = \mathrm{K}(\mathrm{E}^{p})\) , and let \((x_{i})_{i \in \mathrm{I}}\) be a p-basis of E over F (V, §13, No.~1, p.~98). Set \(y_{i} = x_{i} \otimes 1\) and \(z_{i} = x_{i} \otimes 1 - 1 \otimes x_{i}\) . Let \(\Lambda\) be the subset of \(\mathbf{N}^{(I)}\) consisting of the families \((\alpha_{i})_{i \in \mathrm{I}}\) such that \(\alpha_{i} < p\) for every i. Prove that the elements \(\prod_{i \in I} y_{i}^{\alpha_{i}} z_{i}^{\beta_{i}}\) for \((\alpha_{i})\) , \((\beta_{i}) \in \Lambda\) form a basis of the F-vector space \(E \otimes_{F} E\) . Deduce that I is finite (observe that the ring \(E \otimes_{F} E\) is Noetherian).
\item
  Conclude that E has finite degree over K (use Exercise 1, b) of V, §13, p.~170).
\end{enumerate}

¶ 8) Let E be an algebraic extension of a commutative field K. Prove that the ring \(E \otimes_{K} E\) is Noetherian if and only if E has finite degree over K (reduce to the case of a purely inseparable extension using the method of Exercise 6, and then apply Exercise 7).

¶ 9) Let K be a commutative field and A a K-algebra. Suppose that for every Artinian K-algebra B, the ring \(A \otimes_{K} B\) is Artinian. Prove that A is finite-dimensional over K (reduce to the case when A is semisimple by observing that the \(A/\Re(A)\) -modules \(\Re(A)^{k}/\Re(A)^{k+1}\) have finite length; then use Exercises 6 and 7).

\begin{enumerate}
\def\labelenumi{\arabic{enumi})}
\setcounter{enumi}{9}
\tightlist
\item
  Let K be a commutative ring, A a K-algebra, and P an (A,A)-bimodule. An extension of A by P is a triple \((B,i,p)\) , where B is a K-algebra, \(p:B\to A\) a surjective homomorphism of K-algebras, and \(i:P\to B\) an injective K-linear homomorphism with image Ker p, satisfying \(i(p(b)x)=bi(x)\) and \(i(xp(b))=i(x)b\) for \(x\in P\) and \(b\in B\) . An isomorphism from an extension \((B,i,p)\) to an extension \((B',i',p')\) is an isomorphism of K-algebras \(u:B\to B'\) such that \(p'\circ u=p\) and \(u\circ i=i'\) .
\end{enumerate}

\begin{enumerate}
\def\labelenumi{\alph{enumi})}
\item
  Let \((B, i, p)\) be an extension of A by P; the image of i is a two-sided ideal of B with square zero. Conversely, if \(p : B \to A\) is a surjective homomorphism of K-algebras whose kernel P has square zero and i is the canonical injection of P into B, then the triple \((B, i, p)\) is an extension of A by P.
\item
  Let \(i_{0}: P \to A \oplus P\) and \(p_{0}: A \oplus P \to A\) be the canonical mappings. There exists a unique K-algebra structure on \(A \oplus P\) such that \((A \oplus P, i_{0}, p_{0})\) is an extension of A by P. Prove that the automorphism group of this extension can be identified with the group of K-derivations from A to P.
\item
  Let \((B, i, p)\) be an extension of A by P. Prove that the following conditions are equivalent:
\end{enumerate}

\begin{enumerate}
\def\labelenumi{(\roman{enumi})}
\item
  The extension \((\mathrm{B},i,p)\) is isomorphic to the extension \((\mathrm{A}\oplus \mathrm{P},i_0,p_0)\) .
\item
  There exists a homomorphism of K-algebras \(s: \mathrm{A} \to \mathrm{B}\) such that \(p \circ s = \operatorname{Id}_{\mathrm{A}}\) .
\item
  There exists a subalgebra \(S\) of \(B\) such that \(B = S \oplus i(P)\) .
\end{enumerate}

If these conditions are satisfied, then we say that the extension \((B,i,p)\) is trivial. d) Let \((B,i,p)\) and \((B',i',p')\) be extensions of A by P. Let C be the subalgebra of \(B \times B'\) consisting of the elements \((b,b')\) such that \(p(b) = p'(b')\) , let c be the two-sided ideal of C consisting of the elements \((i(x), -i'(x))\) for \(x \in P\) , and let \(B''\) be the algebra C/c and \(\pi : C \to B''\) the canonical homomorphism. Let \(i'' : P \to B''\) and \(p'' : B'' \to A\) be the homomorphisms defined by \(i''(x) = \pi(i(x),0)\) and \(p''(\pi(b,b')) = p(b)\) for \(x \in P\) , \(b \in B\) , \(b' \in B'\) . Prove that \((B'', i'', p'')\) is an extension of A by P; it is called the amalgamated sum of \((B,i,p)\) and \((B',i',p')\) . Prove that we thus define a commutative group law on the set \(\mathrm{Ex}_{\mathrm{K}}(\mathrm{A},\mathrm{P})\) of isomorphism classes of extensions of A by P, whose neutral element is the class of the trivial extension. Moreover, the action of K on \(\mathrm{Ex}_{\mathrm{K}}(\mathrm{A},\mathrm{P})\) given by \(\lambda(\mathrm{B},i,p) = (\mathrm{B},\lambda^{-1}i,p)\) for \(\lambda \in K\) , \(\lambda \neq 0\) defines a K-module structure on \(\mathrm{Ex}_{\mathrm{K}}(\mathrm{A},\mathrm{P})\) .

\begin{enumerate}
\def\labelenumi{\alph{enumi})}
\setcounter{enumi}{4}
\tightlist
\item
  Let \(u: P \to P'\) be a homomorphism of (A, A)-bimodules, and let \((B, i, p)\) be an extension of A by P. Let \(B'\) be the amalgamated sum \(B \oplus_{P} P'\) (II, §1, p.~381, Exercise 5); for \(b \in B\) , \(x' \in P'\) , denote the class of \((b, x')\) in \(B'\) by \([b, x']\) . Let \(i': P' \to B'\) , \(p': B' \to A\) , and \(v: B \to B'\) be the mappings define by
\end{enumerate}

\[
i ^ {\prime} (x ^ {\prime}) = [ 0, x ^ {\prime} ], \qquad p ^ {\prime} ([ b, x ^ {\prime} ]) = p (b), \qquad v (b) = [ b, 0 ].
\]

Prove that there exists a unique algebra structure on \(B'\) such that \((\mathrm{B}', i', p')\) is an extension of A by \(P'\) and v an algebra homomorphism. We say that \((\mathrm{B}', i', p')\) is the extension of A by \(P'\) deduced from \((\mathrm{B}, i, p)\) via u; every extension \((\mathrm{C}, j, q)\) of A by \(P'\) for which there exists an algebra homomorphism \(w: B \to C\) satisfying \(q \circ w = p\) and \(w \circ i = j \circ v\) is isomorphic to \((\mathrm{B}', i', p')\) .

\(f\) ) The mapping \(\mathrm{Hom}_{(\mathrm{A},\mathrm{A})}(\mathrm{P},\mathrm{P}')\times \mathrm{Ex}_\mathrm{K}(\mathrm{A},\mathrm{P})\to \mathrm{Ex}_\mathrm{K}(\mathrm{A},\mathrm{P}')\) that sends a pair \((u,(B,i,p))\) to the class of the extension deduced from (B, \(i,p)\) via \(u\) is K-bilinear. \(g\) ) Let \(f:\mathrm{A}'\to \mathrm{A}\) be a K-algebra homomorphism, and let (B, \(i,p)\) be an extension of A by P. Define likewise the structure of an algebra on the fibered product \(\mathrm{B}'' = \mathrm{B}\times_{\mathrm{A}}\mathrm{A}'\) and an extension \((\mathrm{B}'' , i'', p'')\) of \(\mathrm{A}'\) by P, said to be deduced from (B, \(i,p)\) via \(f\) .

\begin{enumerate}
\def\labelenumi{\alph{enumi})}
\setcounter{enumi}{7}
\tightlist
\item
  Suppose, moreover, that the algebra A is commutative. Prove that the classes of extensions \((B,i,p)\) of A by P such that the algebra B is commutative form a K-submodule of \(\operatorname{Ex}_{\mathrm{K}}(\mathrm{A},\mathrm{P})\) , which we denote by \(\operatorname{Ex}_{\mathrm{K}}^{c}(\mathrm{A},\mathrm{P})\) .
\end{enumerate}

\begin{enumerate}
\def\labelenumi{\arabic{enumi})}
\setcounter{enumi}{10}
\tightlist
\item
  \begin{enumerate}
  \def\labelenumii{\alph{enumii})}
  \tightlist
  \item
    Let f be an element of \(\mathrm{C}^{2}(\mathrm{A},\mathrm{P})\) (VIII, p.~239). Prove that we define the structure of an algebra on the K-module \(A \oplus P\) by setting \((a,x)(a',x') = (aa',ax' + xa' + f(a,a'))\) . The canonical injection \(i: P \to A \oplus P\) and the canonical surjection \(p: A \oplus P \to A\) are ring homomorphisms, and the triple \((\mathrm{A} \oplus \mathrm{P}, i, p)\) is an extension of A by P (observe that we have \(f(a,1) = af(1,1)\) and \(f(1,a) = f(1,1)a\) for every \(a \in A\) ).
  \end{enumerate}
\end{enumerate}

\begin{enumerate}
\def\labelenumi{\alph{enumi})}
\setcounter{enumi}{1}
\item
  Prove that the isomorphism class of this extension depends only on the class of f in \(\mathrm{H}^{2}(\mathrm{A},\mathrm{P})\) . So the construction in a) defines an isomorphism of K-modules from \(\mathrm{H}^{2}(\mathrm{A},\mathrm{P})\) to the K-submodule of \(\mathrm{Ex}_{\mathrm{K}}(\mathrm{A},\mathrm{P})\) consisting of the extensions of A by P that are trivial as extensions of K-modules.
\item
  Suppose that the ring A is commutative; denote by \(\mathrm{C}^{2}(\mathrm{A},\mathrm{P})^{s}\) the K-submodule of \(\mathrm{C}^{2}(\mathrm{A},\mathrm{P})\) consisting of the elements f such that \(f(x,y)=f(y,x)\) for all x,y in A and by \(\mathrm{H}^{2}(\mathrm{A},\mathrm{P})^{s}\) the image of \(\mathrm{C}^{2}(\mathrm{A},\mathrm{P})^{s}\) in \(\mathrm{H}^{2}(\mathrm{A},\mathrm{P})\) . Prove that the isomorphism in b) induces an isomorphism from \(\mathrm{H}^{2}(\mathrm{A},\mathrm{P})^{s}\) to the K-submodule of \(\mathrm{Ex}_{\mathrm{K}}^{c}(\mathrm{A},\mathrm{P})\) consisting of the commutative extensions of A by P that are trivial as extensions of K-modules.
\end{enumerate}

\begin{enumerate}
\def\labelenumi{\arabic{enumi})}
\setcounter{enumi}{11}
\tightlist
\item
  Let A be a commutative ring, J an ideal of A, B the ring A/J, and P a B-module. a) Prove that the exact sequence \(0 \rightarrow J/J^{2} \rightarrow A/J^{2} \rightarrow B \rightarrow 0\) defines an extension of the A-algebra B by \(J/J^{2}\) (Exercise 10).
\end{enumerate}

\begin{enumerate}
\def\labelenumi{\alph{enumi})}
\setcounter{enumi}{1}
\tightlist
\item
  Let P be a B-module. Prove that the homomorphism \(\operatorname{Hom}_{\mathrm{B}}(\mathrm{J}/\mathrm{J}^{2},\mathrm{P}) \to \operatorname{Ex}_{\mathrm{A}}^{c}(\mathrm{B},\mathrm{P})\) that sends u to the extension deduced from the extension in a) via u is an isomorphism.
\end{enumerate}

